% TOP_PROBLEM: {"id": 1246, "title": "Scholz Conjecture", "category_id": 1, "category": {"id": 1, "name": "number_theory", "display_name": "Number Theory", "description": "Properties of integers, prime numbers, Diophantine equations.", "slug": "number-theory", "order_index": 1, "created_at": "2026-07-31T15:26:25.670Z"}, "status": "open"}
% ATTEMPT_STATUS: unresolved
% Research attempt by GPT_6_Astra_Ultra, 1 October 2026.
% Canonical UnsolvedMath ID; original statements and sources preserved.
% FIRST_POSTED: 2026-10-01
% Imported from: unsolvedmath_additional_500.tex; UnsolvedMath ID 1246
% !TeX program = lualatex
% Compile twice: lualatex 1246.tex
% Self-contained: no external bibliography, images, or \input files.
% Numeric section labels and visible numbers are original UnsolvedMath IDs.
\documentclass[11pt,a4paper]{article}
\usepackage[margin=24mm,headheight=14pt]{geometry}
\usepackage{fontspec}
\setmainfont{Latin Modern Roman}
\setsansfont{Latin Modern Sans}
\setmonofont{Latin Modern Mono}
\usepackage{amsmath,amssymb,mathtools}
\usepackage{unicode-math}
\setmathfont{Latin Modern Math}
\usepackage{microtype}
\usepackage{enumitem,xurl,needspace}
\usepackage[dvipsnames]{xcolor}
\usepackage{hyperref}
\hypersetup{unicode=true,colorlinks=true,linkcolor=MidnightBlue,urlcolor=MidnightBlue,
 pdftitle={UnsolvedMath Problem 1246},pdfauthor={GPT 6 Astra Ultra},bookmarksnumbered=true}
\usepackage{bookmark,tocloft,titlesec,fancyhdr}
\setlength{\cftsecnumwidth}{4.2em}
\cftsetpnumwidth{2.5em}
\cftsetrmarg{4em}
\titleformat{\section}{\Large\bfseries\raggedright}{\thesection}{0.7em}{}
\pagestyle{fancy}\fancyhf{}
\fancyhead[L]{\small\sffamily GPT 6 Astra Ultra research attempt}
\fancyhead[R]{\small Original catalogue IDs}
\fancyfoot[C]{\thepage}
\setlength{\parindent}{0pt}
\setlength{\parskip}{5pt plus 1pt}
\setlength{\emergencystretch}{3em}
\setcounter{tocdepth}{1}\setcounter{secnumdepth}{1}
\setlist[itemize]{leftmargin=3em,itemsep=4pt,topsep=4pt,parsep=0pt}
\allowdisplaybreaks
\newcommand{\N}{\mathbb{N}}
\newcommand{\Z}{\mathbb{Z}}
\newcommand{\Q}{\mathbb{Q}}
\newcommand{\R}{\mathbb{R}}
\newcommand{\C}{\mathbb{C}}
\newcommand{\F}{\mathbb{F}}
\newcommand{\A}{\mathbb{A}}
\newcommand{\PP}{\mathbb{P}}
\newcommand{\E}{\mathbb{E}}
\newcommand{\eps}{\varepsilon}
\newcommand{\dd}{\,\mathrm d}
\newcommand{\sref}[2]{\hyperlink{source:#1:#2}{[S#2]}}
\newif\ifshowcatalogue
\showcataloguetrue
\newcommand{\cataloguescope}[1]{\ifshowcatalogue
 \paragraph{Statement recorded in the supplied source.}
 #1\fi}
\begin{document}
% UnsolvedMath ID 1246; source catalogue code NT-039

\setcounter{section}{1245}

\section{Scholz's Addition-Chain Conjecture}\label{1246}

{\small\sffamily UnsolvedMath ID 1246 · Number Theory}

\subsection{Definitions and mathematical statement}

\paragraph{Shared notation.}Unless a section says otherwise, $\N=\{1,2,\ldots\}$,
$\Z$ is the ring of integers, $\Q,\R,\C$ are the rational, real, and complex
fields, $[N]=\{1,\ldots,\lfloor N\rfloor\}$, and $\log$ is the natural
logarithm. For real functions of a parameter tending to infinity,
$f\ll g$ and $f=O(g)$ mean $|f|\leq Cg$ eventually for some fixed $C$;
$f\gg g$ reverses this comparison; $f=o(g)$ means $f/g\to0$;
$f\sim g$ means $f/g\to1$; and $f\asymp g$ means both order bounds.
A subscript on $O$ or $\ll$ specifies allowed constant dependence.
The expression $n^{o(1)}$ in an upper bound means growth smaller than
$n^\epsilon$ for each fixed $\epsilon>0$. Local definitions govern any
exception, finite-field convention, or use of the same symbol for another
object. An asymptotic claim is not automatically an effective algorithm.



An addition chain for $m\geq1$ is $1=a_0<a_1<\cdots<a_\ell=m$ such that each $a_i=a_j+a_k$ for some $j,k<i$, with $j=k$ allowed. Its length is $\ell$, and $\ell(m)$ is the least possible length. The conjecture is $\ell(2^n-1)\leq n-1+\ell(n)$ for every $n\geq1$. \sref{1246}{1} \sref{1246}{2}

\paragraph{Statement recorded in the supplied source.}
Is the shortest addition chain for $2^n - 1$ at most $n - 1$ plus the length of the shortest addition chain for $n$? \sref{1246}{1}

\paragraph{Residual task reported by the archive.}

The embedded review (2026-08-17) records: Prove $l(2^n-1)\leq n-1+l(n)$ for every positive integer n or find a genuine counterexample. \sref{1246}{1}

\subsection{Short English statement}

Can an efficient addition chain for n always be turned into one for two to the nth power minus one with at most n minus one extra steps?

\subsection{Research attempt}

\paragraph{Scope and selected construction.}
The bound involves the shortest unrestricted addition-chain length
$\ell(n)$, not the length of a convenient binary chain. The traditional
star-chain construction associated with Brauer [S2] is useful but has a
specific limitation. Clift's account [S3], inspected 1 October 2026,
discusses that limitation and more elaborate chain classes. The proof
below is self-contained for star chains and is not presented as a proof
that every shortest chain has the required form.

\paragraph{The elementary lower bound on length.}
Every entry in a length-$r$ addition chain is at most $2^r$ at step $r$:
at each new step the sum of two earlier entries is at most twice the
current maximum. Thus
\[
                      \ell(n)\ge\lceil\log_2n\rceil.
\]
The case $n=1$ has length zero and the conjectured bound reads
$\ell(1)\le0$, which holds. For later constructions the number one is
already available and costs no additional step.

\paragraph{Lifting a star chain to a Mersenne-number chain.}
A star chain is an addition chain
$1=a_0<a_1<\cdots<a_r=n$ in which every step is
$a_i=a_{i-1}+a_{j_i}$ with $j_i<i$. Suppose that the values
$M_{a_j}=2^{a_j}-1$ have already been produced for $j<i$, and the
last produced value is $M_{a_{i-1}}$. The identity
\[
 M_{a_i}=2^{a_{j_i}}M_{a_{i-1}}+M_{a_{j_i}}
\]
produces the next desired value using $a_{j_i}$ doublings followed by
one addition of the stored $M_{a_{j_i}}$. Every intermediate value is
strictly larger than the previous one, and every addition uses earlier
available entries, so this is a valid addition chain, not a computation
with uncounted multiplication operations.

The cost is
\[
 \sum_{i=1}^r(a_{j_i}+1)
  =\sum_{i=1}^r(a_i-a_{i-1})+r=n-1+r.
\]
Writing $\ell_*(n)$ for the shortest star-chain length proves
\[
                  \ell(2^n-1)\le n-1+\ell_*(n).
\]
In particular the conjecture holds for every $n$ admitting a shortest
chain that is a star chain. The telescoping identity is the crucial
reason this special construction has the right cost.

\paragraph{Infinite subcases with optimal star chains.}
If $n=2^a$, the repeated-doubling chain has length $a$, which matches
the lower bound and is a star chain. Thus the desired bound holds for
all powers of two. More generally let
$n=2^a+2^b$ with $a>b\ge0$. First build $2^{a-b}$ by $a-b$
doublings, add one, and then double $b$ times. This gives a star chain
of length $a+1$. Since $2^a<n<2^{a+1}$, the lower bound is also
$a+1$. The chain is shortest, so the conjecture holds whenever the
binary expansion of $n$ has at most two nonzero bits.

For arbitrary $n$, the usual left-to-right binary algorithm is a star
chain with length $\lfloor\log_2n\rfloor+s_2(n)-1$, where $s_2(n)$
counts the ones in the binary expansion. Each further bit costs one
doubling, and each further one costs an addition of one. Consequently
the lifting proves the unconditional weaker bound
\[
 \ell(2^n-1)\le n-2+\lfloor\log_2n\rfloor+s_2(n).
\]
This is a bound for every $n$, but replacing its binary-chain length by
$\ell(n)$ is a further step, not an algebraic simplification.

\paragraph{A small worked construction.}
For $n=5$ take the optimal chain $1,2,4,5$. Lifting it gives
\[
                  1,2,3,6,12,15,30,31.
\]
Its seven additions equal $5-1+\ell(5)=4+3$. The stored smaller
Mersenne values are essential: at the exponent-doubling step from two
to four, the value $M_2=3$ is retained and added to $12$ to obtain
$15$. Merely replacing each exponent entry by $2^{a_i}-1$ would give
$1,3,15,31$, which is not an addition chain on its own.

\paragraph{Why the same count fails for arbitrary chains.}
For a general step $a_i=a_j+a_k$ with neither summand necessarily
$a_{i-1}$, the identity
$M_{a_i}=2^{a_k}M_{a_j}+M_{a_k}$ is still valid. But its doublings
start from an older value, and the increments $a_k$ no longer telescope
as $a_i-a_{i-1}$. Some newly needed intermediate values may also be
less than the current maximum, requiring a reordered construction or
additional stored values. This is a cost and dependency problem, not
a failure of the algebraic identity.

Thus a claim that every addition-chain step lifts with total doubling
cost $n-1$ must prove a global schedule or another invariant. Summing
local costs without such a schedule can overcount or omit shared
intermediate work. Clift [S3] gives more advanced examples where a
simple common-path accounting fails; the present argument does not
assume those examples prove the conjecture false, nor does it need
their computational claims to establish the star-chain result.

\paragraph{Executed optimality and construction checks.}
For every $1\le n\le32$, the script finds a shortest unrestricted
chain by iterative-deepening exhaustive search. At each depth it permits
the sum of any two earlier entries, not only star steps; it prunes only
when repeated doubling cannot reach the target in the remaining steps.
Testing all depths from the elementary lower bound up to the first
success certifies optimality in this finite range. Each returned chain
happens to be a star chain. The script lifts it and checks every output
step and the exact cost $n-1+\ell(n)$. It therefore verifies the desired
inequality through $n=32$ by explicit upper-bound certificates, without
claiming to find $\ell(2^n-1)$ exactly.
Files: \texttt{checks.py}, \texttt{results.json} in
\path{_Local/Erdos_problems/UnsolvedMath/attempt_audit_2026_10_01/number_theory/}.

\paragraph{Remaining gap and outcome.}
The construction proves the conjecture on the optimal-star-chain
subclass and, in particular, on the two infinite binary subfamilies
above. What remains is a construction with the target cost for arbitrary
$n$ using its unrestricted optimum, or another argument proving that
bound without lifting a shortest chain. \textbf{Outcome: UNRESOLVED.}
No equality between $\ell_*(n)$ and $\ell(n)$ is assumed in general.

\subsection{Sources}

{\small

\begin{itemize}

\item[\hypertarget{source:1246:1}{S1}] ULAM.AI, \emph{UnsolvedMath}, supplied \texttt{problems.json}, record 1246 (NT-039), statement and background. The embedded literature-review date is 2026-08-17; that is the archive’s review date, not a fresh certification. Dataset landing page: \url{https://huggingface.co/datasets/ulamai/UnsolvedMath}.

\item[\hypertarget{source:1246:2}{S2}] Alfred Brauer, On addition chains, Bulletin of the American Mathematical Society 45 (1939), 736-739. \url{https://projecteuclid.org/journals/bulletin-of-the-american-mathematical-society/volume-45/issue-10/On-addition-chains/bams/1183502523.full}. Bibliographic information imported from the supplied record.

\item[\hypertarget{source:1246:3}{S3}] Neill Clift, The Scholz-Brauer Conjecture. \url{https://additionchains.com/ScholzBrauer.html}. Bibliographic information imported from the supplied record.


\item[E1] Neill Clift, \emph{The Scholz--Brauer Conjecture},
author's computational and structural account.
\url{https://additionchains.com/ScholzBrauer.html}.
Inspected 1 October 2026. Used to distinguish construction restrictions
from the unrestricted conjecture; its large searches were not rerun.

\end{itemize}

}

\label{end:1246}
\end{document}
